\documentclass[11pt]{article}
\usepackage{arxiv}

\title{Pre-Compiled Pipeline Shards for Distributed LLM Inference\\on Intel AI PC Fleets}

\author{
  Tate Berenbaum \\
  Community Labs \\
  \texttt{tb@communitylabs.com}
  \and
  Muthiah Venkatachalam \\
  Intel Corporation \\
  \texttt{muthaiah.venkatachalam@intel.com}
}

\date{April 2026}

\begin{document}

\maketitle

\begin{abstract}
We distribute large language model inference across fleets of commodity Intel AI PCs using pre-compiled pipeline shards composed with mask-based speculative decoding. Per-stage INT4 OpenVINO IR subgraphs are exported via \texttt{torch.jit.trace} with precomputed rotary embeddings and a post-export \texttt{beam\_idx} Gather injection that unlocks the OpenVINO GPU plugin's \texttt{IndirectKVCache} fusion; the resulting shards run at parity with monolithic inference. On top of the shard pipeline we layer two composing multipliers. First, a \emph{mask-based} KV-cache rewind makes speculative decoding practical on stateful OpenVINO models without paged attention---a $\sim 48$~ms \texttt{query\_state}/\texttt{set\_state} round-trip per step is the reason the technique was previously a net loss on Intel iGPUs; masking rejected-draft cache positions via \texttt{attention\_mask}\,=\,0 is bit-exact with physical trim at essentially zero cost and yields 1.35$\times$ mean speedup (1.65$\times$ at 2048 tokens) on a single node. Second, multi-user pipeline micro-batching via independent stateful \texttt{InferRequest}s scales the two-node pipeline 2.03$\times$ in system throughput over the single-stream distributed baseline. Composed end-to-end, a two-node Llama~3.1~8B INT4 pipeline serves two users at 41.25~tok/s aggregate (20.6~tok/s per user)---1.89$\times$ the throughput of monolithic single-user inference on the same hardware---and the advantage grows to 4.04$\times$ over naïve distributed decode under 100~ms/hop simulated WAN latency, where the naïve pipeline falls below the interactive threshold. We also document an undocumented OpenVINO runtime failure mode (\texttt{reset\_state()} on multiple \texttt{InferRequest}s sharing one compiled model) and provide a per-stream \texttt{compile\_model} workaround. Code, raw benchmark logs, and reproduction scripts are released at \texttt{[URL]}.
\end{abstract}

% ============================================================================
\section{Introduction}
\label{sec:intro}
% ============================================================================

Modern AI PCs ship with integrated GPUs, NPUs, and 16--32~GB of unified memory, yet spend most of their time idle. A single Intel AI PC runs Llama~3.1~8B at 22~tokens/second. A fleet of them, coordinated over a standard office network, could serve LLM inference without cloud dependency---at zero marginal cost, with full data locality.

The barrier is software, not hardware. Standard model export tools (\texttt{torch.export}, \texttt{torch.onnx.export}) fail on modern transformer attention due to dynamic control flow in rotary embeddings and KV cache management. Monolithic compiled models cannot be split at layer boundaries without invasive graph surgery. Pipeline parallelism, developed for training throughput~\cite{huang2019gpipe, narayanan2021pipedream}, has seen limited evaluation for autoregressive \emph{inference} on consumer hardware---and when stages are distributed over a real network, each per-token TCP round trip becomes the dominant cost. Speculative decoding~\cite{leviathan2023fast, chen2023accelerating} is the natural answer (amortize network round-trips across multiple tokens emitted per target forward), but the canonical implementations rely on paged attention or dedicated rewind APIs that are not available on stateful OpenVINO models.

We build a distributed-inference stack that closes all three gaps, and the composition produces a surprisingly strong result: a two-node fleet of consumer Intel AI PCs serves two concurrent users at $1.89\times$ the throughput of a single-user monolithic baseline on the same hardware, and under a 100~ms/hop simulated WAN the same configuration stays usable (11.20~tok/s aggregate, 5.6~tok/s per stream) while naïve pipeline-parallel decode falls to 2.77~tok/s (unusable).

We make four contributions:

\begin{enumerate}[leftmargin=*]
  \item \textbf{A per-stage export pipeline that produces INT4 OpenVINO IR shards at monolithic parity.} Starting from \texttt{torch.jit.trace} with precomputed rotary embeddings~(\S\ref{sec:export}) and adding a post-export \texttt{beam\_idx} Parameter + \texttt{Gather(ReadValue, beam\_idx, axis=0)} injection that unlocks the OpenVINO GPU plugin's \texttt{IndirectKVCache} transformation~(\S\ref{sec:shards}), the resulting shards reach within $4\%$ of \texttt{openvino\_genai.LLMPipeline} monolithic inference. Without the \texttt{beam\_idx} injection, the same shards run $12$--$25\%$ slower; the fusion unlock is the non-trivial part of the pipeline, and we document how we found it.

  \item \textbf{Mask-based KV-cache rewind for speculative decoding on stateful models.} Speculative decoding on an OpenVINO stateful target requires rejecting speculative draft tokens from the KV cache. The natural path---\texttt{query\_state()}/\texttt{set\_state()} physical trim---costs $\sim 48$~ms per call on Arc B390 iGPU with a 72-token cache (measured; one trim per step amortizing over the rejected drafts in that step), enough to make speculative decoding a net loss at break-even accept rates. Setting \texttt{attention\_mask[j]\,=\,0} for rejected-draft cache positions is bit-exact with physical trim at essentially zero cost and enables $1.35\times$ mean single-node speedup across eight prompt types, rising to $1.65\times$ at 2048-token generations~(\S\ref{sec:spec}). We document this as Discovery~\#20 in the accompanying repository.

  \item \textbf{Pipeline micro-batching via independent stateful InferRequests.} Each InferRequest carries its own KV cache. Interleaving decode steps from separate user requests across pipeline stages yields $2.03\times$ system-throughput scaling on our $v_5$\_beam Llama~8B shards~(\S\ref{sec:microbatch}). We also document an undocumented OpenVINO runtime failure mode---\texttt{reset\_state()} on multiple \texttt{InferRequest}s sharing one compiled model corrupts shape inference with a specific \texttt{Add} broadcast error---and provide a per-stream \texttt{compile\_model} workaround that recovers the advertised independence semantics.

  \item \textbf{Multiplicative composition of all three into a distributed-inference stack that beats monolithic single-user on throughput and beats naïve distributed decode by $4\times$ under simulated WAN.} The $v_5$\_beam shards $\times$ mask-based spec decode $\times$ 2-stream micro-batching on a two-node fleet reaches $41.25$~tok/s ($1.89\times$ monolithic single-user) at LAN and degrades gracefully with simulated hop latency (measured at $0$, $10$, $50$, $100$~ms/hop)~(\S\ref{sec:distributed}). Naïve distributed decode at $100$~ms/hop yields $2.77$~tok/s; our stack yields $11.20$~tok/s aggregate.
\end{enumerate}

The system is evaluated on three Intel AI PCs (two ASUS Zenbook~S~14 with Lunar Lake, one HP OmniBook~X~16 with Panther Lake) connected over WiFi. Results are reproducible: every table row below is backed by a committed benchmark script and raw log in the companion repository.

% ============================================================================
\section{Background and Related Work}
\label{sec:related}
% ============================================================================

\subsection{Pipeline Parallelism for Inference}

Pipeline parallelism partitions a model's layers across devices. GPipe~\cite{huang2019gpipe} and PipeDream~\cite{narayanan2021pipedream} established this for training, where micro-batches keep all stages busy. Autoregressive inference is different: each token depends on the previous, creating strict sequential dependencies during decode.

Petals~\cite{borzunov2023petals} demonstrated decentralized pipeline inference across volunteer GPUs over the internet, reaching 6~tok/s on Llama~2~70B via DHT-based routing. Parallax~\cite{lin2024parallax} showed pipeline inference on consumer Apple Silicon Macs, reporting 3.1$\times$ latency improvements. MDI-LLM~\cite{wang2025mdillm} proposed recurrent pipeline parallelism for edge devices with up to 74\% throughput gains on four-node systems.

Our work differs in three respects. We pre-compile OpenVINO IR shards with a \texttt{beam\_idx}-Gather graph-surgery pass that unlocks the \texttt{IndirectKVCache} GPU fusion, producing per-stage shards at monolithic parity rather than partitioning models at runtime with graph overhead. We target Intel AI PC integrated GPUs via OpenVINO rather than NVIDIA or Apple hardware. And we introduce mask-based speculative decoding rewind~(\S\ref{sec:spec})---a technique that makes speculative decoding~\cite{leviathan2023fast, chen2023accelerating} pay off on stateful OpenVINO without paged attention---which, stacked with micro-batching, produces the multi-user throughput advantage the paper builds toward.

\subsection{Model Compilation and Export}

The standard path from PyTorch to a compiled inference graph relies on abstract tensor shape propagation (\texttt{torch.export}, \texttt{torch.onnx.export}), which fails on the reshape and transpose operations inside modern rotary position embeddings. OpenVINO's \texttt{ov.convert\_model} inherits these failures when given an \texttt{nn.Module}.

The \texttt{optimum-intel} library exports whole models via \texttt{OVModelForCausalLM} but produces monolithic graphs with no per-layer splitting mechanism. KTransformers~\cite{chen2025ktransformers} compiles individual MoE experts for CPU/GPU hybrid inference but does not handle per-layer-range export for dense models.

Our pipeline uses \texttt{torch.jit.trace} with real tensors---avoiding abstract shape propagation entirely---and precomputes rotary embeddings externally, passing them as flat inputs. This sidesteps the tracing failures while producing per-stage graphs that the OpenVINO compiler optimizes independently.

\subsection{Serving Systems and Throughput Optimization}

vLLM~\cite{kwon2023vllm} introduced PagedAttention for GPU KV cache management. Sarathi-Serve~\cite{agrawal2024sarathi} proposed chunked prefills with stall-free scheduling. TD-Pipe~\cite{zhang2025tdpipe} temporally disaggregates prefill and decode phases, achieving 1.91$\times$ throughput. Splitwise~\cite{patel2024splitwise} splits these phases across specialized hardware.

These systems target data-center GPU clusters. Our micro-batching is simpler: each OpenVINO InferRequest carries independent KV cache state via ReadValue/Assign ops in the compiled graph, so temporal interleaving of requests requires no framework changes and no memory management beyond what OpenVINO already provides.

% ============================================================================
\section{System Design}
\label{sec:design}
% ============================================================================

\subsection{Overview}

Four components: (1)~a per-stage export pipeline producing compiled INT4 OpenVINO IR shards with stateful KV cache, (2)~a TCP activation relay, (3)~stage workers that load and run their assigned shard, and (4)~a coordinator driving the autoregressive generation loop.

Each node loads one shard and exposes a TCP endpoint. The coordinator tokenizes input, runs stage~0 locally, sends the resulting hidden state to the next node, and receives either the next hidden state (intermediate stages) or output logits (final stage). Generated tokens propagate backward to the coordinator.

\subsection{Per-Stage Export Pipeline}
\label{sec:export}

The pipeline converts a HuggingFace transformer into $N$ standalone OpenVINO IR shards, each covering a contiguous layer range plus (for the first) the embedding and (for the last) the output head.

\textbf{Step 1: Selective layer loading.} Decoder layers are loaded from safetensors via a model-agnostic structure mapping (Llama, Mistral, Qwen, Phi, Gemma). Each stage gets only its assigned layers. Memory usage scales with layer count, not model size.

\textbf{Step 2: Attention rewrite.} HuggingFace's \texttt{DynamicCache} (Python list ops) and \texttt{LlamaRotaryEmbedding} (\texttt{torch.autocast} guards) are both untraceable. We replace them with a manual attention implementation taking precomputed cos/sin tensors and explicit KV tensors as inputs. Numerical equivalence verified: max difference $< 5 \times 10^{-7}$ against HuggingFace's native forward.

\textbf{Step 3: Trace and convert.} Each stage is traced with \texttt{torch.jit.trace} using real tensors (batch size~1, representative sequence length). The traced graph converts to OpenVINO IR via \texttt{ov.convert\_model}. KV cache tensors become stateful ReadValue/Assign pairs via \texttt{make\_stateful\_transformation}.

\textbf{Step 4: INT4 compression.} INT4 symmetric weight compression via NNCF~\cite{nncf2024}, group size~128. Each shard is a self-contained IR (XML~+~BIN) taking token IDs or hidden states in, producing hidden states or logits out.

\subsection{TCP Activation Relay}

Activations travel between nodes via raw TCP with a 20-byte header (payload length, dtype, three shape dimensions). KV-cached decode sends \texttt{[1,\,1,\,4096]} float32 tensors: 16~KB per hop. \texttt{TCP\_NODELAY} is set. The protocol is intentionally minimal---designed to be swapped for RDMA, compression, or other transports without touching the rest of the system.

\subsection{Stateful KV-Cached Inference}

Each shard's compiled graph contains ReadValue/Assign ops implementing the KV cache. Prefill processes the full prompt in one pass; decode processes one token per step, reading cached context and appending the new KV pair. The coordinator precomputes rotary cos/sin tensors for the maximum sequence length and slices the appropriate position range per step.

For models with cross-layer KV sharing (Gemma~4 shares K/V between layers via \texttt{DynamicCache}), the shared tensors are transmitted as additional stage outputs/inputs over TCP, with a 4D$\rightarrow$3D reshape for the 3-dimension header protocol.

% ============================================================================
\section{Reaching Monolithic Parity: the beam\_idx Gather Injection}
\label{sec:shards}
% ============================================================================

A naive reading of the export pipeline of \S\ref{sec:export} gives a workable but uncompetitive result: on Llama~3.1~8B INT4, a 1-stage shard (all 32~layers in one graph, same weights as the monolithic model, only the export path differs) runs $13$--$23\%$ slower than \texttt{openvino\_genai.LLMPipeline} at steady state. Splitting into more stages makes it worse. The reason is that the OpenVINO GPU plugin contains a plugin-internal transformation, \texttt{IndirectKVCache}, that rewrites \texttt{ReadValue~$\rightarrow$~Concat~$\rightarrow$~Assign} patterns into a fused \texttt{IndirectSDPA} op---but the transformation is gated on the presence of a specific pattern: each KV \texttt{ReadValue}'s output must flow through a \texttt{Gather} indexed by a \texttt{beam\_idx} Parameter (used by beam search in the monolithic \texttt{LLMPipeline} path). Our initial exports did not include that pattern, so the GPU plugin silently fell back to a generic \texttt{KVCacheFusion} path with measurably lower XVE occupancy.

The fix is post-export graph surgery: after \texttt{apply\_make\_stateful\_transformation}, walk the graph, add a \texttt{beam\_idx: [-1] i32} Parameter, and for each KV \texttt{ReadValue} insert a \texttt{Gather(ReadValue,~beam\_idx,~axis=0)} between the ReadValue output and its downstream consumers. This is exactly what \texttt{optimum-intel}'s internal \texttt{fuse\_cache\_reorder} pass does for monolithic exports, lifted out and applied to our per-stage exports. We call the resulting shards $v_5$\_beam.

\subsection{Results}

\begin{table}[ht]
\centering
\caption{Single-node Llama~3.1~8B INT4 on Panther Lake Arc B390 iGPU, $N\!=\!5$ paired runs, same session, controlled-tokenization 128-token decode. $v_5$\_beam closes the export-to-monolithic gap from $-10$ to within $4\%$ of \texttt{openvino\_genai}.}
\label{tab:shard_parity}
\begin{tabular}{@{}lrrr@{}}
\toprule
Configuration & Tok/s & Range & vs.\ A \\
\midrule
A: mono \texttt{openvino\_genai} (C++ loop)       & 22.96 & 21.71--23.65 & 1.00$\times$ \\
A$'$: mono \texttt{ov.compile\_model} (our Python loop) & 20.61 & 20.09--21.83 & 0.90$\times$ \\
B ($v_5$\_beam 1-stage, our export)               & 22.08 & 21.66--22.63 & 0.96$\times$ \\
B$_{v_3\_fp32}$ (pre-\texttt{beam\_idx})          & 20.05 & 19.74--20.34 & 0.87$\times$ \\
\bottomrule
\end{tabular}
\end{table}

The \texttt{beam\_idx} Gather injection recovers $10\%$ of throughput (from $20.05$ to $22.08$~tok/s) without touching the weights, tokenization, or decode loop. The remaining $4\%$ gap between B ($v_5$\_beam) and A is explained by \texttt{openvino\_genai}'s C++ decode loop vs.\ our Python loop: the A$'$~vs.\ A comparison, which holds the compiled graph fixed and varies only the decode-loop language, shows a $10\%$ Python overhead. After accounting for that, B~is $7\%$ \emph{faster} than A$'$ on the same paired measurement.

\subsection{Splitting into More Stages}

\begin{table}[ht]
\centering
\caption{$v_5$\_beam shards, 2-stage and 3-stage, in-process on Panther Lake iGPU. Same in-process baseline as the spec-decode and progression tables ($A_{\text{spec}}=21.79$, slightly below Table~\ref{tab:shard_parity}'s $A=22.96$ due to a different paired session and run-to-run variance). The parity extends through splitting at the cost of $\sim 1$~tok/s per additional stage from Python activation passing and smaller per-stage XVE problem sizes.}
\label{tab:shard_stages}
\begin{tabular}{@{}lrr@{}}
\toprule
Configuration & Tok/s & vs.\ $A_{\text{spec}}$ (21.79) \\
\midrule
Mono ($A_{\text{spec}}$, reference) & 21.79 & 1.00$\times$ \\
$v_5$\_beam, 1-stage             & 22.08 & 1.01$\times$ \\
$v_5$\_beam, 2-stage (16+16)     & 19.41 & 0.89$\times$ \\
$v_5$\_beam, 3-stage (11+11+10)  & 19.79 & 0.91$\times$ \\
\bottomrule
\end{tabular}
\end{table}

Per-stage splitting costs $\sim$9--11\% on a single machine relative to $A_{\text{spec}}=21.79$ (and 12--15\% relative to Table~\ref{tab:shard_parity}'s tighter $A=22.96$, where $v_5$\_beam 1-stage = 0.96$\times$); the per-stage GEMMs are smaller, reducing XVE occupancy (measured $60.1\%$ on mono vs.\ $43.6\%$ on 3-stage via VTune \texttt{gpu-hotspots}, with identical top-kernel counts), and the Python-level \texttt{hidden\_states.astype(np.float32)} cast between stages is not free. This gap is structural: smaller per-stage problems cannot saturate the GPU the way a single large problem does. We return to how to recover this cost (and then some) via multi-user micro-batching in~\S\ref{sec:microbatch} and via speculative decoding in~\S\ref{sec:spec}.

\subsection{Methodology}

The numbers in Table~\ref{tab:shard_parity} come from one paired session using identical tokenization (chat template applied to both paths so \texttt{input\_ids} are byte-identical), identical \texttt{max\_new\_tokens}~=~128 with \texttt{ignore\_eos=True} (so EOS-driven early termination cannot bias short runs), and identical decode loops except where noted (A vs.\ A$'$ isolates loop language, A$'$ vs.\ B isolates the export path). $N\!=\!5$ runs after 2 warmups. Replication commits, raw JSON, and VTune summary reports are in the companion repository.

% ============================================================================
\section{Speculative Decoding via Mask-Based KV Rewind}
\label{sec:spec}
% ============================================================================

Speculative decoding~\cite{leviathan2023fast, chen2023accelerating} amortizes target-model forward passes over multiple emitted tokens: a smaller draft model proposes $K$ tokens, the target verifies all $K$ in one forward pass, and accepted prefixes are emitted together. On a distributed pipeline this is doubly attractive---it amortizes not only target compute but also network round-trips. On a \emph{single} node, the technique's per-step cost structure determines whether it pays off at all.

\subsection{The Trim Problem}

Standard speculative decoding rejects draft tokens whose predictions do not match the target's. On an OpenVINO stateful model, rejecting tokens means removing their K and V contributions from the KV cache before the next target forward. The natural API for this is \texttt{InferRequest.query\_state()} (returns a list of the per-layer state tensors) followed by assigning a sliced tensor back to each state via \texttt{sv.state~=~ov.Tensor(sliced\_np)}. We measured this on Llama~3.1~8B INT4 on Arc B390:

\begin{table}[ht]
\centering
\caption{Per-call cost of physical KV trim on Arc B390 iGPU with a 72-token cache (32 layers, 64 state tensors). All three variants cluster around $45$--$50$~ms; the cost is device-side state invalidation, not numpy work.}
\label{tab:trim_cost}
\begin{tabular}{@{}lr@{}}
\toprule
\texttt{trim\_kv} variant & ms/call \\
\midrule
\texttt{np.asarray(sv.state.data)} + slice + new Tensor & 49.5 \\
In-place \texttt{Tensor.set\_shape()} & 46.6 \\
\texttt{ov.Tensor} allocate + \texttt{np.asarray(new)[:]} & 49.9 \\
\bottomrule
\end{tabular}
\end{table}

A $K\!=\!3$ speculative step executes one target verify ($\sim 47$~ms memory-bound forward of $K\!+\!1\!=\!4$ tokens), 2--3 draft feeds ($\sim 16$~ms each), and one correction feed ($\sim 16$~ms). At roughly 3 emitted tokens per step, the $\sim 48$~ms trim per step (Table~\ref{tab:trim_cost}) is enough to push the whole technique below break-even.

\subsection{Mask-Based Rewind}

Our alternative avoids state manipulation entirely. The OpenVINO stateful model takes \texttt{attention\_mask} as a per-forward input of shape $[\mathrm{batch},\,\mathrm{past\_len}+\mathrm{input\_len}]$, where \texttt{mask[j]=0} means ``ignore cache position $j$ in attention.'' If rejected-draft tokens remain physically in the cache but their positions are masked out on every subsequent forward, the attention result is identical to a physically-trimmed cache. Future tokens' key and value tensors append at the next cache index, with RoPE computed from the caller-supplied \texttt{position\_ids}; as long as \texttt{position\_ids} tracks the \emph{logical} sequence length (the pre-rewind count) rather than the physical cache length, the rotary rotations are coherent.

We verified this is bit-exact with physical trim. Running both paths side-by-side on the same target--prompt--draft--correction sequence (script \texttt{mask\_trim\_test.py} in the companion repo), the maximum absolute difference in post-correction logits was $0.0000$ (fp32 equality to machine precision).

The cost is CPU-only: each forward reconstructs the $[1, \mathrm{past\_len}+\mathrm{input\_len}]$ mask by concatenating a tracked \texttt{valid\_mask} prefix with ones for the new tokens. Measured: $<\,1\%$ of per-step wall time at $K\!=\!3$. The memory cost is that rejected-draft K and V entries persist in cache, inflating cache size over generations. At 2048-token generation with $K\!=\!3$, cache bloat reaches $1.02\times$ logical length (most rejected-draft overhead is amortized out by the live-token count), so no periodic compaction is needed for realistic single-turn lengths.

\subsection{Single-Node Results}

\begin{table}[ht]
\centering
\caption{Speculative decoding on Llama~3.1~8B INT4 (target) + Llama~3.2~1B INT4 (draft) on Arc B390 iGPU, greedy decode, bit-exact output vs.\ target-only greedy. Eight diverse prompts, 128-token decode, $K\!=\!3$, mean $\pm$ stddev.}
\label{tab:spec_lan}
\begin{tabular}{@{}lrrr@{}}
\toprule
Prompt category & Baseline (tok/s) & Spec (tok/s) & Speedup \\
\midrule
short-factual    & 22.08 & 28.75 & 1.30$\times$ \\
reasoning        & 20.60 & 27.27 & 1.32$\times$ \\
code-completion  & 19.98 & 32.71 & 1.64$\times$ \\
list-enumeration & 21.10 & 28.54 & 1.35$\times$ \\
creative         & 22.10 & 24.53 & 1.11$\times$ \\
technical-expl   & 21.98 & 27.12 & 1.23$\times$ \\
chat-assistant   & 21.96 & 30.73 & 1.40$\times$ \\
translation      & 22.15 & 31.93 & 1.44$\times$ \\
\midrule
\textbf{Mean}    & \textbf{21.49} & \textbf{28.95} & \textbf{1.35$\times$ $\pm$ 0.15} \\
\bottomrule
\end{tabular}
\end{table}

Speedup varies from $1.11\times$ (creative writing, 49.7\% accept rate) to $1.64\times$ (code completion, 93.1\% accept rate). On content that is easier for the draft to predict (structured output: code, translation, chat boilerplate), the speedup climbs toward the theoretical $1\!+\!K\!\cdot\!p_\text{accept}$ ceiling; on content where the draft diverges early, it falls to the acceptance floor. Every prompt produces output that is bit-exact identical to the target's own greedy decode on the first ten generated tokens (the companion repo's smoke test verifies full-sequence equality).

Generation length amplifies the effect: across 128, 256, 1024, and 2048-token runs on the \emph{same} prompt, acceptance rises from $70.7\%$ to $97.5\%$ as the model enters its repetition regime, and speedup at $K\!=\!3$ climbs from $1.35\times$ to $1.65\times$.

\subsection{Compose-with-Shards}

Because the rewind lives at the Python wrapper layer---not inside the compiled OV graph---the same \texttt{MaskedReq} abstraction composes with both the monolithic target and our $v_5$\_beam 3-stage sharded target. Measured in a single paired session (\texttt{bench\_spec\_matrix.py}):

\begin{table}[ht]
\centering
\caption{Speculative decoding stacks cleanly with sharding. Single-session paired measurement, $K\!=\!3$, 128-token decode.}
\label{tab:spec_shard}
\begin{tabular}{@{}lrrr@{}}
\toprule
Target & Tok/s & vs.\ mono-A & vs.\ own baseline \\
\midrule
Mono (A) & 21.79 & 1.00$\times$ & --- \\
3-stage $v_5$\_beam shard (B) & 19.79 & 0.91$\times$ & --- \\
Mono + spec $K\!=\!3$ & 28.28 & 1.30$\times$ & 1.30$\times$ \\
3-stage shard + spec $K\!=\!3$ & 26.28 & 1.21$\times$ & 1.33$\times$ \\
\bottomrule
\end{tabular}
\end{table}

The spec-decode multiplier is essentially the same ($1.30\times$ vs.\ $1.33\times$) on both paths. Output is bit-exact: the mono-plus-spec output matches the mono-only output, and the shard-plus-spec output matches the shard-only output.

% ============================================================================
\section{Distributed Pipeline Evaluation}
\label{sec:distributed}
% ============================================================================

\subsection{Testbed}

Three Intel AI PCs on the same WiFi network (802.11ax):

\begin{itemize}[leftmargin=*]
  \item \textbf{node-00} (\emph{beta}, ASUS Zenbook S 14): Core Ultra 7 258V (Lunar Lake), Arc 140V iGPU, 32~GB LPDDR5X-8533.
  \item \textbf{node-01} (\emph{charlie}, ASUS Zenbook S 14): Core Ultra 7 258V (Lunar Lake), Arc 140V iGPU, 32~GB LPDDR5X-8533.
  \item \textbf{node-02} (\emph{alpha}, HP OmniBook X 16): Core Ultra X7 358H (Panther Lake), Arc B390 iGPU, 32~GB DDR5.
\end{itemize}

All run Windows~11, Python~3.11, OpenVINO~2026.1.0, GPU inference. Raw TCP round-trip for 16~KB payloads: 6.5~ms on the 802.11ax LAN (measured separately).

\subsection{Progressive Optimization}

\begin{table}[ht]
\centering
\caption{Llama~3.1~8B INT4 distributed throughput on the 2-node \emph{alpha} + \emph{charlie} $v_5$\_beam pipeline, building up from the single-stream distributed baseline. System throughput for multi-user is aggregated across concurrent streams. All configurations produce bit-exact output relative to target-only greedy. The full-stack number ($41.25$) is from one paired session of \texttt{mini\_coord\_spec\_mbatch.py} (commit \texttt{ed84f27}); the WAN sweep in Table~\ref{tab:wan} reports $41.01$ at the LAN datapoint of a separate session (commit \texttt{6ab4da3}). The $0.24$~tok/s gap is run-to-run variance well within the typical $\pm 0.3$ stddev observed across sessions.}
\label{tab:progression}
\begin{tabular}{@{}lrrr@{}}
\toprule
Configuration & Tok/s & vs.\ mono-$A_{\text{spec}}$ & vs.\ single-stream dist. \\
\midrule
Mono single-node reference ($A_{\text{spec}}$) & 21.79 & 1.00$\times$ & --- \\
2-stage $v_5$\_beam, single stream (\emph{alpha}+\emph{charlie}) & 14.51 & 0.67$\times$ & 1.00$\times$ \\
\,+ 2-stream micro-batching (\S\ref{sec:microbatch}) & 29.46 & 1.35$\times$ & 2.03$\times$ \\
\,+ mask-based spec decode $K\!=\!3$ (\S\ref{sec:spec}) & \textbf{41.25} & \textbf{1.89$\times$} & \textbf{2.84$\times$} \\
\bottomrule
\end{tabular}
\end{table}

Starting from the 2-node $v_5$\_beam pipeline at $14.51$~tok/s ($23$~ms slower per token than the $21.79$~tok/s single-node mono reference; of that $23$~ms, $\sim 6$~ms is one TCP round-trip and $\sim 15$~ms is Python wrapper plus OpenVINO dispatch overhead, see Table~\ref{tab:breakdown}), two composable multipliers stack multiplicatively:

\textbf{Micro-batching ($2.03\times$ over single-stream distributed).} Two independent \texttt{InferRequest}s per shard let the coordinator interleave stage-0 compute for request~B while stage~1 is busy with request~A. Our $v_5$\_beam micro-batch multiplier ($29.46 / 14.51 = 2.03\times$) reflects the smaller per-stage compute window that \texttt{IndirectKVCache} fusion leaves behind---more idle time for the other stream to fill (\S\ref{sec:microbatch_compare} compares against an earlier internal export iteration). (This $2.03\times$ is the mbatch multiplier; not to be confused with the $2.03\times$ system multiplier of full-stack-K=2 over mono-A in \S\ref{sec:wan}, which is a different ratio.)

\textbf{Mask-based spec decode ($1.40\times$ on top of mbatch).} Each of the two streams runs its own speculative decoding loop (\S\ref{sec:spec}). The $K\!+\!1$-token target verify batches four tokens per TCP round-trip and amortizes the per-hop latency. Applied atop the two-stream pipeline, this yields a further $\sim 1.40\times$ multiplier, for a system-throughput total of $41.25$~tok/s: $1.89\times$ the monolithic single-user reference on the same hardware, two users served concurrently.

Composition is empirical, not a projection: $14.51 \times 2.030 \times 1.400 = 41.24$, matching the measured $41.25$ to one significant digit; all three multipliers were measured in the same paired session (\texttt{scripts/mini\_coord\_spec\_mbatch.py}, commit \texttt{ed84f27}).

\subsection{Per-Token Breakdown}

\begin{table}[ht]
\centering
\caption{Where the time goes: 2-stage $v_5$\_beam Llama~3.1~8B distributed pipeline on the \emph{alpha}+\emph{charlie} 802.11ax WiFi LAN, GPU--GPU, single stream (no spec decode, no micro-batching). Compute numbers averaged across N=3 timed runs.}
\label{tab:breakdown}
\begin{tabular}{@{}lrr@{}}
\toprule
Component & Time & \% of token \\
\midrule
Stage~0 compute (coordinator, \emph{alpha} Panther Lake) & 23~ms & 34\% \\
Stage~1 compute (worker, \emph{charlie} Lunar Lake) & 25~ms & 36\% \\
TCP round trip (single hop, 802.11ax) & 6~ms & 9\% \\
Python activation passing + mask construction & 15~ms & 21\% \\
\midrule
\textbf{Total} & \textbf{69~ms} & $\rightarrow$ \textbf{14.51~tok/s} \\
\bottomrule
\end{tabular}
\end{table}

An important subtlety: our initial instrumentation reported ``70\% network time,'' but that metric conflated actual TCP latency with the time spent waiting for the remote worker to compute. Dedicated TCP benchmarking measured the true network cost at $\sim 6$~ms per round trip ($9\%$ of per-token time; the round-figure $6.5$~ms cited elsewhere is the same number rounded differently across sessions). The rest is compute and Python overhead (the latter is OpenVINO's C++ dispatch and state management, not user Python---a dedicated profile run (\texttt{bench\_feed\_overhead.py}) finds user Python is $0.6\%$ of wall time).

\subsection{Activation Compression}

\begin{table}[ht]
\centering
\caption{Activation compression, 2-stage distributed Llama~3.1~8B. On LAN, bandwidth is not the bottleneck---compression cannot help.}
\label{tab:compression}
\begin{tabular}{@{}lrrrl@{}}
\toprule
Compression & Payload/hop & Tok/s & Decode & Quality \\
\midrule
None (FP32) & 16~KB & 13.57 & 65--68~ms & Correct \\
FP16 & 8~KB & 12.71 & 66--67~ms & Correct \\
INT8 & 4~KB & 15.10 & 63--68~ms & Corrupted \\
\bottomrule
\end{tabular}
\end{table}

At gigabit WiFi, 16~KB transmits in under $0.2$~ms. The $\sim 6.5$~ms per-hop latency is TCP kernel scheduling and round-trip time, not bandwidth. Halving the payload gains nothing. INT8 symmetric quantization of the 4096-dim hidden state corrupts generation---``What is the capital of France?'' yields ``1.\ Paris 2.\ London 3.\ Berlin 4.\ Rome.'' Per-channel or group quantization might preserve quality but was not tested. Activation compression is only relevant under bandwidth-constrained WAN conditions.

\subsection{WAN Latency Sensitivity}
\label{sec:wan}

We inject one-way latency on the worker node's \texttt{recv()} and \texttt{send()} boundaries (so added per-target-forward RTT~$\approx 2 \times$ injected), sweeping across LAN-to-intercontinental values. Unlike traditional \texttt{tc/netem} WAN emulation, this approach is deterministic but lacks real TCP effects like slow-start or jitter---the reported speedups should be considered a lower bound on the real WAN advantage (real TCP congestion control amortizes multiple in-flight small messages; the sleep-sim serializes them).

\begin{table}[ht]
\centering
\caption{Full-stack Llama~3.1~8B throughput vs.\ simulated hop latency, 2-stage $v_5$\_beam pipeline, 2 concurrent streams with spec decode $K\!=\!3$. Naïve single-stream distributed decode (no mbatch, no spec) shown as baseline. At $100$~ms/hop the naïve path falls below the interactive floor while the full stack remains usable.}
\label{tab:wan}
\begin{tabular}{@{}lrrrr@{}}
\toprule
Hop latency (ms) & Naïve baseline & Full stack (agg) & Per-stream & Full / naïve \\
\midrule
0 (LAN)  & 14.51 & 41.01 & 20.5 & 2.83$\times$ \\
10       &  9.89 & 37.35 & 18.7 & 3.78$\times$ \\
50       &  4.72 & 18.25 &  9.1 & 3.87$\times$ \\
100      &  2.77 & 11.20 &  5.6 & \textbf{4.04$\times$} \\
\bottomrule
\end{tabular}
\end{table}

The advantage over naïve distributed decode \emph{grows} with hop latency. At LAN the full-stack multiplier is $2.83\times$ naïve; at $100$~ms/hop it is $4.04\times$. The full stack composes \emph{two} amortizations: micro-batching folds two streams into one pipeline (an $\approx 2\times$ ceiling for balanced stages), and speculative decoding's $K+1$-token target verify amortizes per-hop latency by $1+p_\text{accept}\cdot K$. For our measured $K\!=\!3$ with $p_\text{accept}\!=\!0.707$, the spec-only ceiling is $3.12\times$ over the spec-less single-stream baseline; combined with the mbatch $\approx 2\times$, the full-stack ceiling against naïve is $\approx 6.2\times$. The observed $4.04\times$ at $100$~ms/hop falls between the mbatch-alone and full-ceiling bounds, consistent with neither amortization being saturated at this hop latency.

\paragraph{Real-WAN cross-check.} Because the sleep-sim approach in Table~\ref{tab:wan} only models the latency dimension (not the TCP cwnd dynamics that govern multi-segment messages), we ran a parallel sweep with two release-time-queue TCP latency proxies on the alpha node forwarding to charlie/beta with per-byte one-way delay. The proxy queues each chunk with a release timestamp and a separate sender thread forwards it on time, so multiple chunks can be in flight (matching tc-netem behavior, not per-chunk serialization). Run on the 3-stage Llama target-only path at the time of this revision, with the same prompt and 64-token decode budget per run:

\begin{table}[ht]
\centering
\caption{Real-WAN sweep with on-alpha release-time-queue TCP proxies (3-stage Llama~3.1~8B v5\_beam, target-only single stream, 2 runs per latency, 64-token decode). Compared to the sleep-sim's single-hop 2-stage column in Table~\ref{tab:wan}, this row pays \emph{two} hops per token (alpha~$\to$~charlie, alpha~$\to$~beta) and gets no spec/mbatch amortization. The gap to the sleep-sim 2-stage figures comes from the topology change (2~hops/token vs.\ 1) and the absence of speculative decoding---not from the proxy method itself.}
\label{tab:real_wan}
\begin{tabular}{@{}lr@{}}
\toprule
Hop latency (ms, one-way) & Tok/s (3-stage target-only, real WAN) \\
\midrule
0   (LAN)  & 12.53 \\
10         &  8.18 \\
50         &  3.28 \\
100        &  2.01 \\
\bottomrule
\end{tabular}
\end{table}

Adding hop latency on the 3-stage target-only path costs more wall-clock per token than the sleep-sim suggests, because each token requires four full TCP round-trips with logits responses approaching $500$~KB --- TCP cwnd takes several RTTs to ramp on a cold flow, and persistent connections only partially amortize the warm-cwnd penalty. This is not a flaw in either sleep-sim (which captures the latency dimension correctly for stream-level behavior) or the queue proxy (which captures real TCP dynamics); it is evidence that for a paper-quality comparison, both methods should be reported. We retain Table~\ref{tab:wan}'s sleep-sim numbers as the headline because that pipeline includes spec~$+$~mbatch (matching the full-stack story); the queue-proxy numbers above isolate the 3-stage TCP overhead.

Choosing the right $K$ per latency regime matters. A separate sweep on the 3-stage $v_5$\_beam target-only pipeline (single stream, no mbatch; \texttt{scripts/bench\_spec\_wan\_K.py}) finds that LAN prefers mid $K$ ($K\!=\!7$ wins, $K\!=\!5$ close) where compute dominates and excess draft feeds waste cycles, while $100$~ms/hop prefers large $K$ ($K\!=\!10$) where batching enough tokens per round-trip dominates.

\begin{table}[ht]
\centering
\caption{$K$-sweep on the 3-stage $v_5$\_beam pipeline (target-only single stream, no mbatch), measured 2026-04-24 in one paired session of \texttt{bench\_spec\_wan\_K.py}. Baselines (no spec): LAN $21.55$~tok/s, $50$~ms/hop $4.77$, $100$~ms/hop $2.77$. Optimal $K$ shifts from $7$ at LAN (compute-bound, large drafts waste cycles when compute is cheap) to $10$ at $100$~ms/hop (network-bound, packing more tokens per round-trip dominates).}
\label{tab:k_sweep}
\begin{tabular}{@{}lrrrrrrr@{}}
\toprule
$K$ & accept & LAN tok/s & LAN $\times$ & 50~ms tok/s & 50~ms $\times$ & 100~ms tok/s & 100~ms $\times$ \\
\midrule
2  & 84.2\% & 28.63 & 1.33$\times$ & 10.27 & 2.15$\times$ &  6.47 & 2.33$\times$ \\
3  & 70.7\% & 26.77 & 1.24$\times$ & 11.54 & 2.42$\times$ &  7.31 & 2.64$\times$ \\
5  & 67.1\% & 30.42 & 1.41$\times$ & 14.05 & 2.94$\times$ &  9.28 & 3.34$\times$ \\
7  & 58.9\% & \textbf{30.96} & \textbf{1.44$\times$} & 15.27 & 3.20$\times$ & 10.51 & 3.79$\times$ \\
10 & 52.5\% & 28.04 & 1.30$\times$ & \textbf{16.04} & \textbf{3.36$\times$} & \textbf{11.30} & \textbf{4.07$\times$} \\
\bottomrule
\end{tabular}
\end{table}

The LAN and WAN regimes of Table~\ref{tab:k_sweep} give two distinct operating points:

\textit{LAN / cloud-adjacent.} Compute dominates. $K\!=\!7$ peaks at $30.96$~tok/s ($1.44\times$ over the single-stream LAN baseline of $21.55$); $K\!=\!5$ trails closely at $30.42$ ($1.41\times$). Small-$K$ benefits from high acceptance but the extra draft feeds at small $K$ cost more than they save when compute is already cheap; large-$K$ wastes drafts. The sweet spot is moderate $K$ ($5$--$7$) that balances acceptance with draft cost; precise winner shifts within that band across runs by $\sim 3\%$.

\textit{Regional-to-intercontinental WAN.} Network dominates. $K\!=\!10$ at $50$~ms/hop reaches $16.04$~tok/s ($3.36\times$ over naïve $4.77$ at the same hop latency); at $100$~ms/hop it reaches $11.30$~tok/s ($4.07\times$ over naïve $2.77$). This is the headline shard-specific advantage measured on the target-only 3-stage path. Note that Table~\ref{tab:wan} reports a similar $4.04\times$ at $100$~ms/hop using a different configuration (2-stream mbatch + spec $K\!=\!3$, full stack); both regimes hit roughly the same multiplier through different mechanisms because at high latency the network term dominates the per-token wall time and the multipliers approach each other.

\subsection{Long Generation and Stability}

Decode latency remains flat over extended generations. On the single-node monolithic baseline, throughput actually \emph{improves} with generation length---speculative decoding's acceptance rate rises as the model enters more predictable continuation patterns:

\begin{itemize}[leftmargin=*]
  \item $128$ tokens: $1.35\times$ spec speedup, $70.7\%$ acceptance at $K\!=\!3$ (mean across 8 prompts; \texttt{bench\_spec\_v7\_masked.py}).
  \item $512$ tokens: $1.55\times$ speedup, $90.6\%$ acceptance (\texttt{bench\_spec\_long\_gen.py}).
  \item $1{,}024$ tokens: $1.59\times$ speedup, $95.1\%$ acceptance.
  \item $2{,}048$ tokens: $1.65\times$ speedup, $97.5\%$ acceptance, cache bloat ratio $1.02\times$ (the mask-based rewind leaves rejected-draft K/V in physical cache, but rejected-draft overhead is per-step constant while the live-token count grows linearly---bloat ratio converges to 1).
\end{itemize}

On the 2-stage distributed pipeline (single-stream, no spec), earlier measurements on an extended-prompt workload showed no per-token degradation across $200$ tokens ($15.95$~tok/s), $1000$ tokens ($15.55$~tok/s), and 10 consecutive prompts ($14.46$~tok/s aggregate, $0.72$~QPS, with $36$~ms KV-cache reset between prompts). Those numbers predate $v_5$\_beam; we have not re-measured sustained long-generation on the full stack.

\subsection{Gemma 4 E2B: A Second Architecture}
\label{sec:gemma}

The distributed pipeline extends to Gemma~4~E2B (5.1B, 35 layers, FP32 due to PLI quantization sensitivity, see \S\ref{sec:discussion}). Gemma~4 also exercises a feature Llama doesn't have: layers~15--34 share K/V projections with earlier layers (the weights are untrained placeholders, and inference reads from source layers' cache via \texttt{DynamicCache}). In the distributed pipeline, this means L13/L14 KV tensors must cross the stage boundary. We transmit them as additional non-stateful outputs from stage~0, received as inputs at stage~1, with a 4D$\rightarrow$3D reshape for the TCP header protocol.

The Gemma path also forced a separate fix on OpenVINO 2026.1: HuggingFace's \texttt{Gemma3nTextRotaryEmbedding} computes cos/sin under \texttt{torch.autocast(enabled=False)} and casts to \texttt{x.dtype} at the end, which \texttt{torch.jit.trace} bakes into the IR as a mixed-precision \texttt{opset1::MatMul}. OpenVINO~2026.0 silently auto-promoted the operand types; 2026.1 validates element types strictly and refuses to compile (rotary block has \texttt{f16[?,256,1] x f32[?,1,?]} on GPU). We replace HF's rotary submodule with a custom \texttt{GemmaTracedRotaryEmbedding} that holds \texttt{inv\_freq} as a buffer, computes everything in FP32, and casts to target dtype once at the end. Two instances handle Gemma~4's two layer types: \texttt{sliding\_attention} (default rope, \texttt{head\_dim=256}, $\theta=10\text{k}$) and \texttt{full\_attention} (proportional rope, \texttt{head\_dim=512}, $\theta=1\text{M}$, \texttt{partial\_rotary\_factor=0.25} --- only the first 25\% of dims rotated, the rest pad with zero \texttt{inv\_freq}). The rotary-fixed shards (we call them ``v2'') compile at default GPU precision and recover full speed; the previously-reported numbers were preserved by re-export rather than the \texttt{INFERENCE\_PRECISION\_HINT="f32"} workaround (which costs ${-}10\%$ on 2-stage, ${-}60\%$ on 1-stage).

We also apply the same post-export \texttt{beam\_idx} Gather injection from \S\ref{sec:shards} to Gemma's stage~0 (stage~1's KV-shared layers have no \texttt{ReadValue} ops to inject into, so it is left as-is). The injected variant is reported below as ``v2\_beam.''

\begin{table}[ht]
\centering
\caption{Gemma~4~E2B 2026-04-24 (FP32, OV~2026.1.0, alpha Panther Lake B390 iGPU + charlie Lunar Lake Arc 140V iGPU). The ``v2'' column is the rotary-fixed export; ``v2\_beam'' adds post-hoc \texttt{beam\_idx} Gather injection on stage~0 to unlock the GPU plugin's \texttt{IndirectKVCache} fusion. Output byte-correct on every config.}
\label{tab:gemma_dist}
\begin{tabular}{@{}lrrr@{}}
\toprule
Configuration & 2026-04 paper & v2 (this work) & v2\_beam (this work) \\
\midrule
1-stage single-node (GPU) & 13.3 & 13.69~$\pm$~0.44 & \textbf{14.13~$\pm$~0.21} \\
2-stage localhost GPU$\rightarrow$GPU & 12.1 & 12.32~$\pm$~0.32 & \textbf{13.13~$\pm$~0.14} \\
2-stage multi-node GPU$\rightarrow$GPU & 8.12 & 10.49~$\pm$~0.20 & 10.66~$\pm$~0.27 \\
2-stage multi-node CPU$\rightarrow$GPU & 7.16 & \textbf{9.77~$\pm$~0.04} & --- \\
2-stream micro-batch (multi-node, agg) & 13.51 & \textbf{16.39~$\pm$~0.07} & 16.15~$\pm$~0.03 \\
\bottomrule
\end{tabular}
\end{table}

The single-stream multi-node figure of $\mathbf{10.49}$~tok/s on v2 (and $10.66$ with \texttt{beam\_idx} fusion) is $29\%$ better than the 2026-04 paper's $8.12$. The improvement composes two effects: the v2 export removes a precision-hint penalty present in any compile attempted on OV~2026.1+, and the Lunar Lake-class iGPU on charlie has matured with the OV~2026.1 release. The $10.66$~tok/s v2\_beam delta over v2 ($+1.6\%$) is small in distributed mode because per-token wall time is dominated by the 8--16~KB hidden-state and 2-pair cross-KV TCP round-trips, not on-device compute --- consistent with the 2-stage in-process result where v2\_beam beats v2 by $6.6\%$ (compute-bound), and weakly with the 1-stage in-process where the gain is $3.2\%$ (the 1-stage path is also compute-bound but has only one beam-injected stage rather than one-of-two).

The CPU stage~0 path was initially blocked by an OpenVINO 2026.1 CPU-plugin shape-inference bug at the empty-state KV-cache concat: after \texttt{reset\_state()}, the \texttt{ReadValue} carries shape \texttt{f32[0,1,0,256]}, and the subsequent \texttt{Concat(ReadValue, RoPE:f32[1,1,16,256])} fails validation because dim~0 differs between the two operands. The GPU plugin tolerates the empty leading dim; the CPU plugin rejects it. The workaround is to explicitly initialize each state to shape \texttt{[1, num\_kv\_heads, 0, head\_dim]} via \texttt{infer\_request.query\_state()[i].state = ov.Tensor(np.zeros(...))} immediately after \texttt{reset\_state()}; this is documented in the export script's verification pass and now applied uniformly in \texttt{gemma\_2s\_coord.py}. With the workaround, CPU$\rightarrow$GPU multi-node runs at $\mathbf{9.77}$~tok/s, $36\%$ above the 2026-04 paper's $7.16$.

The 2-stream micro-batch row uses a Gemma-specific multi-stream coordinator and worker (per-stream \texttt{compile\_model} on both sides as the documented workaround for the \texttt{reset\_state()} multi-\texttt{InferRequest} bug in \S\ref{sec:ov_bug}). Aggregate throughput is $\mathbf{16.39}$~tok/s on v2 ($1.56\times$ over single-stream) and $16.15$~tok/s on v2\_beam ($1.51\times$). Both exceed the paper's $13.51$ by ${\approx}21\%$. The mbatch \emph{ratio} ($1.56\times$/$1.51\times$) is below the paper's $1.66\times$ because the v2 single-stream baseline is itself higher --- there is less stage-idle time to fill on a $10.49$ baseline than on an $8.12$ one. v2\_beam's mbatch is slightly under v2 because under 2-stream contention both compiled stage~0s share GPU compute, so the per-stream wall time loses any margin from the \texttt{IndirectKVCache} fusion. Output between the two concurrent streams is byte-identical (same prompt, deterministic generation), confirming that the cross-stream KV isolation is correct.

% ============================================================================
\section{Multi-User Throughput via Micro-Batching}
\label{sec:microbatch}
% ============================================================================

A 2-stage pipeline has an inherent inefficiency for single requests: while stage~1 computes, stage~0 sits idle. With balanced 25~ms stages, utilization is roughly 50\%.

\subsection{Approach}

Each OpenVINO InferRequest \emph{should}, per the Model API documentation, maintain independent KV cache state. Creating two InferRequests per shard---each with its own cache---lets us interleave two user requests:

\begin{enumerate}[leftmargin=*]
  \item Stage~1 processes request~A's token; simultaneously, stage~0 processes request~B's token.
  \item Stage~1 moves to request~B; stage~0 advances request~A.
  \item Repeat. The pipeline alternates, filling bubbles.
\end{enumerate}

\subsection{An undocumented OpenVINO constraint}
\label{sec:ov_bug}

In practice, OpenVINO~2026.1 and~2026.2 contain an undocumented failure mode: calling \texttt{reset\_state()} on one \texttt{InferRequest} among several created from the \emph{same} \texttt{CompiledModel} causes subsequent \texttt{infer()} calls on any of the sibling requests to fail with:

\begin{verbatim}
Check 'TRShape::broadcast_merge_into(...)' failed at
  src/core/shape_inference/include/eltwise_shape_inference.hpp:28:
While validating node 'opset1::Add Add_NNN () -> ()':
Argument shapes are inconsistent.
\end{verbatim}

The specific \texttt{Add\_NNN} node varies between runs and models; consistently, it is an \texttt{opset1::Add} whose broadcast merge resolves one argument shape to the empty tuple \texttt{()}. The failure does not appear on the first prefill, but reliably occurs after at least one full run-and-reset cycle. We have been unable to find an existing report of this issue in OpenVINO's issue tracker.

\textbf{Workaround:} call \texttt{compile\_model()} separately per stream so each InferRequest has its own fully-independent compiled graph. The cost is extra compile time (seconds per additional stream on Arc iGPU) and extra GPU memory (${\sim}2$~GB per additional Llama~8B half-stage). The resulting configuration recovers the advertised independence semantics, and is what the $41.25$~tok/s full-stack result depends on.

\subsection{Results}

\begin{table}[ht]
\centering
\caption{Micro-batching: gain depends on stage asymmetry \emph{and} on per-stage compute size. The $v_5$\_beam measurements (alpha+charlie, Llama~3.1~8B INT4) scale $2.03\times$, close to the theoretical ceiling; the smaller per-stage GEMMs of $v_5$\_beam (post-\texttt{IndirectKVCache} fusion) leave more stage-idle time for the other stream to fill than the external-rotary export comparison row.}
\label{tab:microbatch}
\label{sec:microbatch_compare}
\begin{tabular}{@{}llrrr@{}}
\toprule
Model & Export / year & Single-stream & 2-stream & Gain \\
\midrule
Llama 3.1 8B & external-rotary v1 (2026-04) & 14.1 & 19.4 & 1.38$\times$ \\
Llama 3.1 8B & $v_5$\_beam (2026-04, this work) & 14.51 & \textbf{29.46} & \textbf{2.03$\times$} \\
Gemma 4 E2B & external-rotary v1 (2026-04) & 8.12 & 13.51 & 1.66$\times$ \\
\bottomrule
\end{tabular}
\end{table}

The $v_5$\_beam Llama result is the relevant number for the full-stack story (it composes with \S\ref{sec:spec}'s speculative decoding to $41.25$~tok/s system throughput). The $1.66\times$ Gemma~4~E2B micro-batch ratio in this row uses the older external-rotary export ($8.12 \to 13.51$~tok/s); the v2 and v2\_beam shards (Table~\ref{tab:gemma_dist}) reach $1.56\times$ and $1.51\times$ respectively. The slightly lower ratio on the new exports is consistent: both numerator (single-stream) and denominator (2-stream aggregate) move up, but the single-stream baseline moves up by more, leaving less stage-idle time for the second stream to absorb. Absolute aggregate is higher: $16.39$~tok/s on v2 vs the table's $13.51$.

The theoretical 2-stream maximum is $2\times$ for balanced stages. Our measured $2.03\times$ on $v_5$\_beam exceeds this slightly---we attribute it to the GPU runtime coalescing some small-GEMM dispatches when both streams are active, though we have not confirmed this with microbenchmarks.

The $v_5$\_beam shards close the monolithic-parity gap \emph{and} produce smaller per-stage compute windows for micro-batching to fill---a dual benefit of the \texttt{beam\_idx} Gather injection introduced in~\S\ref{sec:shards}.

% ============================================================================
\section{Discussion}
\label{sec:discussion}
% ============================================================================

\subsection{When Is a Fleet Better Than Cloud?}

Three conditions: the organization already owns AI PCs (sunk cost), data must stay on-premises (privacy/compliance), and the latency budget accommodates $\sim$50--70~ms tokens (interactive chat). Under these conditions, a 2-node fleet running the full stack (sharded $+$ micro-batched $+$ speculative) serves two concurrent users of Llama~3.1~8B at $41.25$~tok/s aggregate ($20.6$~tok/s per user) with zero marginal cost and full data locality.

A cloud-hosted Llama~3.1~8B typically achieves $30$--$60$~tok/s but incurs per-token API costs and data residency exposure. The fleet trades peak throughput for cost and privacy---and at $100$~ms/hop cross-continental WAN, the fleet-with-speculative-decoding remains interactive ($11$~tok/s aggregate, $5.6$~tok/s per stream) where a naïve distributed pipeline falls to $2.77$~tok/s (unusable).

\subsection{Scaling Projections}

Llama~3.1~70B (35~GB INT4) would need roughly 7 nodes at 32~GB each. Network overhead scales linearly with hops ($\sim$6.5~ms per WiFi hop), so a 7-stage pipeline adds $\sim$39~ms---still interactive if per-stage compute stays around $\sim$25~ms. Micro-batching can interleave up to $\min(N_\text{users}, N_\text{stages})$ requests, approaching full utilization. Speculative decoding's advantage scales favorably with the number of network hops since it amortizes more round-trips per emitted token, so the full stack's benefit grows with pipeline depth---the observed $4.04\times$ speedup over naïve at $100$~ms/hop with two nodes would, by the same round-trip-amortization argument, be larger at seven nodes.

\subsection{Quantization Challenges}

Gemma~4's Per-Layer Input (PLI) embeddings are a dense $262\text{K} \times 8960$ matrix---not a standard lookup table. INT4 and INT8 quantization of this matrix produces garbage output regardless of group size or scheme. FP32 is the only working precision, inflating stage~0 to 7~GB (exceeding Lunar Lake GPU's $\sim$4~GB allocation limit). A mixed-precision export path (FP32 embedding, INT4 decoder layers) would solve this but requires NNCF to support per-subgraph precision control.

\subsection{Limitations}

The system assumes a trusted, reliable network. There is no fault tolerance, authentication, or encryption. WiFi latency varies by $\pm$2~ms (averaging out over many tokens). All nodes must run identical software stacks (same OpenVINO, same Python). Mixed-device pipelines (GPU~+~CPU or GPU~+~NPU) showed poor results: NPU rejects ReadValue ops entirely, and CPU stages on the same SoC contend for memory bandwidth with the GPU stage.

\textbf{3-stage distributed: full-stack composition deferred.} The 3-node topology (alpha + charlie + beta) runs cleanly on the all-Ethernet LAN: three back-to-back 128-token decodes complete with no abort at $14.42 / 14.44 / 14.47$~tok/s ($\approx$2-stage parity, single stream, no spec, no mbatch; output byte-correct), and a $154$~req/s sustained alpha-charlie smoke test (identical traffic shape) confirms the LAN handles the load with margin to spare. We have not yet measured the full stack (3-stage target~$+$~spec~$+$~mbatch) because the spec~$+$~mbatch coordinator is currently 2-stage-only (\S\ref{sec:microbatch}); extending it to 3-stage is straightforward but deferred. (One Wi-Fi-only configuration of the 3-stage path triggered \texttt{WinError 10053} retransmit timeouts; on Ethernet, no aborts occur. Because Wi-Fi is not part of the production target topology, we do not analyze it further here.)

\textbf{WAN characterization (sleep-sim and queue-proxy, both reported).} The headline sweep in Table~\ref{tab:wan} uses \texttt{time.sleep()} on the worker's recv/send boundaries to inject deterministic one-way delay per target forward. Table~\ref{tab:real_wan} adds a release-time-queue TCP-proxy sweep on the 3-stage Llama target-only path to capture real TCP dynamics (per-byte one-way delay, multi-segment in-flight). The two methods differ: sleep-sim captures latency for the full stack with spec~$+$~mbatch (matching the topology of \S\ref{sec:fullstack}); the queue-proxy captures real per-RTT congestion-window behavior on a different topology (3-stage target-only, no spec, no mbatch). We report both rather than picking one --- a real-WAN measurement of the spec~$+$~mbatch path on a multi-machine testbed remains future work.

\subsection{Negative Results}

Optimization attempts that failed, documented here because the failure modes are informative:

\textbf{Early exit.} We applied the final model's RMSNorm~+~lm\_head to layer~15 outputs, hoping high-confidence tokens could skip the second half of the network. The lm\_head never produced confident predictions at layer~15 (max softmax probability stayed below 50\%), and the $800$~MB matrix multiply added $80$~ms overhead per token. The fundamental issue: the lm\_head is calibrated for layer-31 representations, not layer-15. A purpose-trained lightweight exit head might work but requires calibration data we did not collect.

\textbf{Asymmetric shard planning.} A 12/20 layer split (vs.\ the default 16/16) yielded 2\% improvement---within measurement noise. The bottleneck is network round-trip cost, not compute imbalance between stages.

\textbf{GPU~+~NPU heterogeneity.} The Panther Lake NPU (50~TOPS) rejects ReadValue operations, making it incompatible with stateful KV-cached shards. Non-stateful NPU inference is possible but requires full-sequence recomputation each token, negating any NPU speed advantage.

\textbf{\texttt{KV\_CACHE\_PRECISION} and \texttt{INFERENCE\_PRECISION\_HINT} on Arc iGPU.} We tested all documented OpenVINO GPU precision hints to see if INT8 KV cache would accelerate the memory-bound decode path. On Arc B390, the INT8 KV configuration runs $4\%$ \emph{slower} than default ($20.86$ vs.\ $21.74$~tok/s on monolithic Llama~3.1~8B INT4)---dequantization overhead exceeds the memory bandwidth savings at this sequence length on this GPU. Output is bit-exact across all hint values. Per-stage asymmetric precision would not help for the same reason (the GPU is already memory-bound-limited, not compute-bound).

\textbf{Async / threaded overlap of draft and target.} At LAN, trying to overlap draft drafting on Python threads with target verify (using \texttt{SO\_KEEPALIVE}, \texttt{TCP\_NODELAY}, and GIL-releasing \texttt{infer()} calls) yields no measurable gain: $11.08$ vs.\ $11.05$~tok/s at $K\!=\!10$, $100$~ms/hop simulated. The ``all $K$ drafts accepted'' case where async helps has probability $p^K \approx 0.2\%$ at $K\!=\!10$, so simple speculation of next-step drafts during the target's network wait almost always wastes the speculative work. Tree speculation covering multiple accept-count branches would work but requires a 4D attention mask our OV export does not currently support.

\textbf{C++ port of the spec-decode loop.} We considered rewriting \texttt{spec\_decode\_greedy()} in C++ to reduce Python overhead. A dedicated profile (\texttt{bench\_feed\_overhead.py}) finds Python is $0.6\%$ of wall time ($27$~ms out of $4521$~ms for $K\!=\!3$, 128 tokens); OpenVINO's \texttt{req.infer()} accounts for $99.3\%$. The effort would move nothing. The measured $22\%$ Bad-Speculation rate in VTune P-core analysis is OpenVINO's own runtime dispatch/wait loop, not user Python.

% ============================================================================
\section{Conclusion}
\label{sec:conclusion}
% ============================================================================

A fleet of commodity Intel AI PCs can serve multi-user LLM inference at interactive speeds \emph{above} single-user monolithic throughput on the same hardware, and can do so across network links where naïve pipeline parallelism is unusable. The two-node full-stack Llama~3.1~8B configuration---$v_5$\_beam-exported shards, mask-based speculative decoding, two-stream micro-batching---reaches $41.25$~tok/s system throughput, $1.89\times$ the monolithic single-user baseline, and $4.04\times$ naïve distributed decode under a simulated $100$~ms/hop WAN.

The three composing techniques each address a distinct bottleneck and are individually validated with bit-exact output vs.\ their respective baselines:

\begin{itemize}[leftmargin=*]
\item The \texttt{beam\_idx} Gather injection unlocks the OpenVINO GPU plugin's \texttt{IndirectKVCache} fusion for per-stage exports, producing shards at monolithic parity~(\S\ref{sec:shards}).
\item Mask-based KV-cache rewind sidesteps the $\sim 48$~ms-per-call \texttt{query\_state}/\texttt{set\_state} round-trip that would otherwise make speculative decoding a net loss on stateful OpenVINO~(\S\ref{sec:spec}).
\item Independent-state micro-batching fills the stage-idle windows that $v_5$\_beam fusion leaves smaller---at the cost of a per-stream \texttt{compile\_model} to work around an undocumented OpenVINO \texttt{reset\_state} interaction~(\S\ref{sec:ov_bug}, \S\ref{sec:microbatch})---yielding $2.03\times$ scaling at the 2-stream operating point.
\end{itemize}

Code, export scripts, all raw benchmark logs, and VTune reports are available at \texttt{[URL]}.

% ============================================================================
% References
% ============================================================================
\bibliographystyle{plainnat}
\bibliography{references}

\end{document}
